% !TEX program = xelatex
% AI-effects 프로젝트 — Anthropic Economic Index(AEI) 데이터를 1차 실증 자료로
% 직접 활용한 연구 9편만을 대상으로, 각 연구가 (1) AI 노출도(exposure), (2) AI 사용량(usage)을
% 어떻게 정의하고 논의의 중심을 어디에 두는지, 그리고 이전 세대의 노출도·채택 연구를
% 어떻게 계승·발전시켰는지를 정리한 작업 문헌(report/).
% 컴파일: xelatex -> biber -> xelatex -> xelatex (kotex, biblatex/biber 필요)
\documentclass[11pt]{article}

\usepackage{fontspec}
\usepackage{kotex}
\usepackage[a4paper,margin=2.2cm]{geometry}
\usepackage{longtable}
\usepackage{booktabs}
\usepackage{array}
\usepackage{amsmath}
\usepackage{xcolor}
\usepackage{hyperref}
\usepackage{enumitem}
\usepackage{titlesec}
\usepackage[backend=biber,style=authoryear,maxcitenames=2,uniquelist=false,sorting=nyt]{biblatex}
\addbibresource{../../reference/references.bib}

\hypersetup{colorlinks=true,linkcolor=blue!50!black,citecolor=blue!50!black,urlcolor=blue!50!black}
\newcolumntype{L}[1]{>{\raggedright\arraybackslash}p{#1}}
\setlength{\LTleft}{0pt}
\setlength{\LTright}{0pt}
\renewcommand{\arraystretch}{1.15}
\setlength{\tabcolsep}{4pt}

\title{Anthropic Economic Index 기반 실증연구:\\[1mm]
AI 노출도와 AI 사용량의 정의, 계승, 발전\\[2mm]
\large --- 노출(exposure) 중심 논의에서 사용(usage) 중심 논의로, 그리고 그 결합 ---}
\author{AI-effects 프로젝트}
\date{\today}

\begin{document}
\maketitle

\begin{abstract}
\noindent 본 문헌은 \textbf{Anthropic Economic Index(AEI)}의 실측 Claude 사용 데이터를 핵심 자료로 삼는 실증연구 9편—AEI 공식 시리즈 5편(\citealp{handa-2025-which-economic-tasks-are}[R1]; \citealp{appel-2025-uneven-geographic-and-enterprise-ai}[V3]; \citealp{appel-2026-economic-primitives}[V4]; \citealp{massenkoff-2026-learning-curves}[V5]; \citealp{massenkoff-2026-cadences}[V6])과 동일 데이터를 응용한 후속 연구 4편(\citealp{tamkin-2025-estimating-ai-productivity-gains-from}; \citealp{massenkoff-2026-labor-market-impacts-of-ai}; \citealp{fan-2026-what-countries-use-ai-and}; \citealp{fan-2026-aggregate-gains-from-ai})—만을 대상으로, 각 연구가 (1) \textbf{AI 노출도(exposure)}와 (2) \textbf{AI 사용량(usage)}을 각각 어떻게 정의하며, 논의를 어느 쪽 중심으로 전개하는지를 정리한다. 나아가 이 9편이 이전 세대의 (1) 이론적 노출도 연구(\citealp{felten-2021-occupational-industry-and-geographic-exposure}; \citealp{felten-2023-how-will-language-modelers-like-chatgpt}; \citealp{webb-2020-impact-artificial-intelligence-labor}; \citealp{eloundou-2023-gpts-are-gpts-an-early})와 (2) 채택·적용 연구(기업·근로자 서베이형 adoption 지표)를 각각 어떻게 계승·발전시켰는지를 계보적으로 논증한다. 핵심 결론은 다음과 같다: AEI 계열 9편 중 7편은 노출도를 별도 변수로 다루지 않는 순수 \textbf{사용량 중심(usage-centric)} 연구이며, 오직 \citet{massenkoff-2026-labor-market-impacts-of-ai}만이 Eloundou $\beta$(이론노출)를 실사용 데이터로 게이팅한 \textbf{결합형(hybrid) 지표}("관측노출")를 계산적으로 구축하고, \citet{massenkoff-2026-cadences}는 별도로 설문 기반 "보고노출/기대노출"을 실사용 패턴과 연결하는 \textbf{또 다른 방식의 결합}을 시도한다. 나머지는 이론적 노출 개념을 방법론적으로 재사용하지 않고, 오히려 그 예측과 실측치의 괴리를 보여주는 대비축으로만 인용한다.
\end{abstract}

\tableofcontents

% ============================================================
\section{서론}
% ============================================================

\texttt{CLAUDE.md}가 규정하는 이 저장소의 핵심 목표는 AI exposure(잠재적·이론적 노출)와 AI usage(실현된 실제 사용) 데이터를 결합해, 노출만으로는 포착하지 못하는 실제 채택·활용 효과를 분석하는 것이다. 기존 report들은 이 목표를 여러 각도에서 다뤄왔다—\texttt{01\_AI영향\_실증연구\_정리.tex}는 노출도·채택·실사용 세 갈래 선행연구 지형을 개관했고, \texttt{02\_AI노출도\_지표\_비교.tex}는 노출도 지표 8종의 산출공식을 심화 비교했으며, \texttt{06\_AI채택\_실사용\_데이터방법론.tex}는 노출도를 제외한 채택·실사용 실증연구 28편을 데이터 출처·방법론 축으로 정리했다.

본 문헌은 이 계열에서 한 걸음 더 들어가, \textbf{Anthropic Economic Index(AEI)라는 단일 데이터 계열을 활용한 실증연구만}을 대상으로 삼는다. AEI는 Claude.ai 대화·1P API 트랜스크립트를 프라이버시 보존 분류기(Clio)로 O*NET 과업에 매핑해 발행하는 Anthropic 자체 연구보고서 시리즈로, 2025년 2월 첫 보고서(R1) 이후 2026년 6월까지 약 매 분기 새 웨이브를 발행하며 방법론을 확장해 왔다. 이 시리즈가 특히 주목할 가치가 있는 이유는, \texttt{06\_AI채택\_실사용\_데이터방법론.tex}가 지적했듯 실사용 데이터 기반 연구가 아직 얕은 역사(2023년 이후)와 특정 서비스 이용자 국한이라는 이중의 한계를 안고 있음에도, AEI 시리즈 \textbf{내부}에서 노출도 개념을 사용량 측정과 결합하려는 시도(\citealp{massenkoff-2026-labor-market-impacts-of-ai}의 관측노출; \citealp{massenkoff-2026-cadences}의 보고노출-실사용 연계)가 본 프로젝트가 지향하는 "exposure+usage 결합" 목표와 정확히 같은 방향으로 이미 진행되고 있기 때문이다.

\subsection{대상 문헌 선정 기준과 범위}

정리 대상은 \textbf{AEI 데이터(Claude.ai 대화 로그 또는 1P API 트랜스크립트)를 1차 실증 자료로 직접 분석한 논문 9편}으로 한정한다. \texttt{reference/notes/} 전체에서 "Anthropic Economic Index" 또는 "Claude.ai" 데이터 사용 여부를 기준으로 검색한 결과, 다음 9편이 이 기준을 충족한다(발행순).

\begin{enumerate}[label=(\arabic*)]
  \item \citet{handa-2025-which-economic-tasks-are} --- AEI 최초 보고서(R1)
  \item \citet{appel-2025-uneven-geographic-and-enterprise-ai} --- AEI 3번째 보고서(V3)
  \item \citet{tamkin-2025-estimating-ai-productivity-gains-from} --- AEI 데이터를 이용한 생산성 추정 응용연구
  \item \citet{appel-2026-economic-primitives} --- AEI 4번째 보고서(V4, 경제원시지표)
  \item \citet{massenkoff-2026-labor-market-impacts-of-ai} --- AEI 데이터+Eloundou $\beta$ 결합 응용연구
  \item \citet{massenkoff-2026-learning-curves} --- AEI 5번째 보고서(V5, 학습곡선)
  \item \citet{fan-2026-what-countries-use-ai-and} --- AEI 데이터를 이용한 국가간 확산 응용연구
  \item \citet{massenkoff-2026-cadences} --- AEI 6번째 보고서(V6, cadences)
  \item \citet{fan-2026-aggregate-gains-from-ai} --- AEI 데이터를 이용한 IMF 응용연구
\end{enumerate}

이 중 1·2·4·6·8은 Anthropic이 직접 발행하는 "AEI 보고서 시리즈" 본편(R1, V3, V4, V5, V6)이고, 3·5·7·9는 AEI가 공개한 데이터(허깅페이스 공개 데이터셋 포함)를 활용해 Anthropic 내부 연구자 또는 외부 기관(IMF)이 별도로 수행한 응용연구다. 반면 \texttt{06\_AI채택\_실사용\_데이터방법론.tex}에서 이미 다룬 \citet{bick-2026-what-work-does-generative}—서베이 기반 채택지수를 주자료로 삼되 \citet{handa-2025-which-economic-tasks-are}의 AEI 과업점유율 데이터를 비교·검증용으로 2차 활용하는 연구—는 AEI가 1차 자료가 아니므로 본 문헌의 핵심 9편에서는 제외하고, \S\ref{sec:bick}에서 별도로 그 함의만 다룬다. 한국 문헌 중 AEI를 인용하는 사례(예: 오삼일·오영식 2026; 장지연 외 2024)는 모두 \citet{handa-2025-which-economic-tasks-are}의 자동화-증강 비율(43\%/57\%)을 배경 서술로 인용할 뿐 AEI 데이터를 직접 분석하지 않으므로 제외한다.

\subsection{AEI 시리즈 데이터 계보 개관}

\begin{longtable}{L{3.0cm}L{2.6cm}L{4.0cm}L{6.0cm}}
\caption{분석 대상 문헌 9편 개관}\label{tab:overview}\\
\toprule
\textbf{문헌} & \textbf{발행일(AEI 웨이브)} & \textbf{표본 데이터} & \textbf{핵심 기여}\\
\midrule
\endfirsthead
\multicolumn{4}{l}{\small (표 \thetable\ 계속)}\\
\toprule
\textbf{문헌} & \textbf{발행일(AEI 웨이브)} & \textbf{표본 데이터} & \textbf{핵심 기여}\\
\midrule
\endhead
\bottomrule
\endfoot

\citet{handa-2025-which-economic-tasks-are} &
2025.2(R1, 자료 2024.12--2025.1) &
Claude.ai Free/Pro 대화 100만$\sim$400만건 &
Clio로 대화를 O*NET 과업에 최초 매핑, 자동화(43\%)/증강(57\%) 분류\\

\citet{appel-2025-uneven-geographic-and-enterprise-ai} &
2025.9.15(V3, 자료 2025.8) &
Claude.ai 대화 100만건+1P API 100만건 &
지리적 사용지수(AUI) 최초 정의, 소비자-기업(API) 사용 대비\\

\citet{tamkin-2025-estimating-ai-productivity-gains-from} &
2025.11.5(V3.5 응용) &
Claude.ai 대화 10만건 &
과업당 시간절감률 추정, Hulten 정리로 생산성 집계(+1.8\%p)\\

\citet{appel-2026-economic-primitives} &
2026.1.15(V4) &
Claude.ai+1P API 각 100만건(2025.11) &
5개 "경제원시지표"(복잡성$\cdot$숙련$\cdot$용도$\cdot$자율성$\cdot$성공) 도입\\

\citet{massenkoff-2026-labor-market-impacts-of-ai} &
2026.3.5(V3--V4 응용) &
AEI Claude 사용데이터(2025.8$\cdot$11)+Eloundou $\beta$ &
"관측노출(observed exposure)" 지수 구축, CPS 실업률 DID\\

\citet{massenkoff-2026-learning-curves} &
2026.3.24(V5) &
Claude.ai+1P API 각 100만건(2026.2) &
모델선택(Opus 사용)-임금 관계, tenure별 "학습곡선"\\

\citet{fan-2026-what-countries-use-ai-and} &
2026.4(V5 응용) &
AEI Claude.ai 대화 80만여건(100개국+) &
사용강도(PPML)$\cdot$확산폭(HHI) 분해, AI-diffusion space\\

\citet{massenkoff-2026-cadences} &
2026.6.26(V6) &
Claude 대화(2026.4--6)+AEI Survey(N$\approx$9,700) &
산출물(artifact) 분류, 보고노출-실사용 연계 최초 시도\\

\citet{fan-2026-aggregate-gains-from-ai} &
2026.7(R1--V5 응용) &
AEI 5개 웨이브(2025.1--2026.2) &
노동비용등가(LCE)$\cdot$AI집중지수(ACI) 구축\\
\end{longtable}

표 \ref{tab:overview}가 보여주듯 이 문헌군은 단일 저자집단(Anthropic 소속 연구자, 특히 Tamkin·McCrory·Massenkoff·Appel이 여러 편에 걸쳐 중복 참여)이 약 16개월(2025.2--2026.7) 동안 거의 매 분기 데이터·방법론을 갱신하며 축적한 하나의 누적적 연구 프로그램이며, 외부 응용연구(Fan 2편)조차 같은 자료를 그대로 가져다 쓰는 밀접한 계보를 형성한다.

% ============================================================
\section{AI 노출도(exposure)의 정의: AEI 계열 연구는 이를 어떻게 다루는가}
\label{sec:exposure}
% ============================================================

\subsection{선행 노출도 연구 계보(요약)}

\texttt{02\_AI노출도\_지표\_비교.tex}가 상세히 다룬 대로, AEI 이전의 "AI 노출도" 개념은 \textbf{직업·과업이 AI로 자동화·보완될 잠재력을 사전적으로 추정}하는 지표들로 구성된다.

\begin{itemize}[leftmargin=*]
  \item \citet{felten-2021-occupational-industry-and-geographic-exposure}의 AIOE: 크라우드소싱된 인간 능력(ability) 등급과 AI 응용 진전도(AI Progress)를 곱해 직업단위로 집계한 지수. \citet{felten-2023-how-will-language-modelers-like-chatgpt}는 이를 언어모델 특화판으로 갱신.
  \item \citet{webb-2020-impact-artificial-intelligence-labor}: 특허 텍스트와 직업 과업기술서의 어휘적 중복도로 산출한 노출지수.
  \item \citet{eloundou-2023-gpts-are-gpts-an-early}: GPT-4를 평가자로 삼아 과업별 노출등급(E0--E3, $\beta/\zeta/\varphi$)을 산출하는 LLM 기반 노출 루브릭—AEI 계열과 가장 직접적으로 연결되는 선행연구다.
  \item \citet{hampole-2025-artificial-intelligence-and-the-labor}: 기업$\times$과업 임베딩으로 노출을 대체효과($m$)와 집중도($C$)로 분해.
\end{itemize}

이들의 공통점은 \textbf{실제 AI가 그 과업에 쓰이고 있는지를 관측하지 않고}, "쓰일 수 있는가"라는 잠재력만을 전문가 평가$\cdot$크라우드소싱$\cdot$LLM 판정으로 추정한다는 점이다.

\subsection{AEI 계열 연구별 노출도 취급 방식}

\begin{longtable}{L{3.0cm}L{2.0cm}L{4.0cm}L{6.6cm}}
\caption{9편의 AI 노출도(exposure) 취급 방식}\label{tab:exposure}\\
\toprule
\textbf{문헌} & \textbf{변수로 다루는가} & \textbf{다룬다면 정의} & \textbf{선행 노출연구와의 관계}\\
\midrule
\endfirsthead
\multicolumn{4}{l}{\small (표 \thetable\ 계속)}\\
\toprule
\textbf{문헌} & \textbf{변수로 다루는가} & \textbf{다룬다면 정의} & \textbf{선행 노출연구와의 관계}\\
\midrule
\endhead
\bottomrule
\endfoot

\citet{handa-2025-which-economic-tasks-are} & 아니오 & --- &
Webb(2020)$\cdot$Eloundou et al.(2023)의 예측을 실측치와 대비하는 서론상의 대조군으로만 인용(예: Webb의 임금 90분위 정점 예측 vs 실측 중상위 임금 정점)\\

\citet{appel-2025-uneven-geographic-and-enterprise-ai} & 아니오 & --- & 언급 없음(순수 사용량 기술)\\

\citet{tamkin-2025-estimating-ai-productivity-gains-from} & 아니오 & --- &
\citet{acemoglu-2024-the-simple-macroeconomics-of-ai}의 과업기반 거시모형(노출 개념의 집계 틀)을 계승하되, 노출도 지표 자체는 입력하지 않음\\

\citet{appel-2026-economic-primitives} & 아니오 & --- &
"실효 AI 커버리지"가 개념적으로 노출-사용 결합과 유사하나, 이론적 노출등급을 전혀 사용하지 않고 관측된 성공률$\cdot$시간비중만으로 산출\\

\textbf{\citet{massenkoff-2026-labor-market-impacts-of-ai}} & \textbf{예} &
Eloundou et al.(2023)의 $\beta$(1=LLM단독 충분, 0.5=도구필요, 0=불가)를 \textbf{그대로 입력변수로 사용} &
직접 계승: $\beta\geq 0.5$인 과업만 "이론상 노출"로 인정한 뒤 실사용 데이터로 게이팅\\

\citet{massenkoff-2026-learning-curves} & 아니오 & --- & 언급 없음\\

\citet{fan-2026-what-countries-use-ai-and} & 아니오 & --- &
\citet{felten-2021-occupational-industry-and-geographic-exposure}$\cdot$\citet{eloundou-2023-gpts-are-gpts-an-early}를 "직업-과업 노출지수" 계열로 명시적으로 대비, 자사 강도$\cdot$확산폭 지표를 그 대안으로 위치시킴\\

\textbf{\citet{massenkoff-2026-cadences}} & \textbf{예(설문 기반)} &
응답자가 "AI가 오늘/12개월 후 할 수 있는 업무 비중"에 스스로 응답한 \textbf{보고노출(reported exposure)$\cdot$기대노출(anticipated exposure)} &
Felten$\cdot$Eloundou류의 "전문가/LLM이 판정하는 노출"과 달리 \textbf{사용자 본인이 판정하는 주관적 노출}이라는 제3의 조작화. 관측노출$\cdot$이론노출과의 상관을 직접 비교\\

\citet{fan-2026-aggregate-gains-from-ai} & 아니오 & --- &
\citet{felten-2021-occupational-industry-and-geographic-exposure}$\cdot$\citet{eloundou-2023-gpts-are-gpts-an-early}를 "잠재적 노출" 계열로 서론에서 명시적으로 대비, LCE$\cdot$ACI를 그 대안으로 제시\\
\end{longtable}

표 \ref{tab:exposure}가 보여주듯, \textbf{9편 중 7편은 노출도를 아예 변수로 도입하지 않는다.} 이들은 오히려 노출도 개념을 "우리 연구가 극복하는 선행 패러다임"으로 인용하는 수사적 대비축으로만 활용한다—\citet{handa-2025-which-economic-tasks-are}가 Webb$\cdot$Eloundou의 예측이 실제 사용 패턴과 어긋남을 보이는 것, \citet{fan-2026-what-countries-use-ai-and}와 \citet{fan-2026-aggregate-gains-from-ai}가 "노출은 잠재력일 뿐 실현이 아니다"라는 동일한 논거를 국가 단위로 반복하는 것이 대표적이다. 노출도를 계산적으로 재사용하는 것은 \citet{massenkoff-2026-labor-market-impacts-of-ai} 한 편뿐이며, \citet{massenkoff-2026-cadences}은 이론적 노출이 아니라 \textbf{설문에 기반한 주관적 노출 인식}이라는 별개의 개념을 사용량과 연결한다.

% ============================================================
\section{AI 사용량(usage)의 정의: AEI 계열 연구의 조작화}
\label{sec:usage}
% ============================================================

\subsection{선행 채택$\cdot$적용 연구와의 근본적 차이}

\texttt{06\_AI채택\_실사용\_데이터방법론.tex} \S2가 정리했듯, AEI 이전의 "AI 채택(adoption)" 측정은 압도적으로 \textbf{자기보고 서베이}(정부 기업패널조사, 근로자 서베이)에 의존해 왔다—\citet{bick-2026-mind-the-gap-ai-adoption}의 미국-유럽 근로자 서베이, 한국 기업활동조사$\cdot$사업체패널조사(WPS)의 도입 여부 더미 등. 이들은 "AI를 도입했는가/사용하는가"라는 \textbf{이분법적$\cdot$자기평가적} 질문에 응답자가 답하는 구조이며, 응답자마다 "도입"의 기준이 다를 수 있다는 근본적 한계를 안는다.

AEI 계열이 이 지형에 도입한 방법론적 전환은, 자기보고를 전혀 거치지 않고 \textbf{실제 대화 로그 자체를 관측}해 어떤 과업에 AI가 실제로 쓰였는지를 분류기(Clio)로 직접 판정한다는 점이다. \citet{handa-2025-which-economic-tasks-are}가 스스로 명시하듯, 이는 "기존의 이론적 노출(exposure)지수... 와 대비되는 실측(observed usage) 접근"일 뿐 아니라, 암묵적으로는 자기보고형 채택 서베이와도 대비되는 접근이다—응답자가 "나는 AI를 업무에 쓴다"고 답하는 대신, 그가 실제로 무엇을 입력했고 AI가 무엇을 출력했는지가 직접 분류된다.

\subsection{각 문헌의 사용량 정의와 측정: 방법론적 심화의 계보}

AEI 계열 9편은 "사용량"을 다음과 같이 \textbf{점진적으로 심화}시켜 왔다. 이는 단일 개념의 반복 측정이 아니라, 매 웨이브가 이전 웨이브의 측정 한계를 지적하며 새 차원을 추가하는 누적적 발전이다.

\begin{longtable}{L{2.8cm}L{2.8cm}L{4.8cm}L{5.0cm}}
\caption{9편의 AI 사용량(usage) 정의와 측정, 그리고 계승 관계}\label{tab:usage}\\
\toprule
\textbf{문헌} & \textbf{핵심 지표} & \textbf{조작화 방식} & \textbf{직전 연구 대비 심화 지점}\\
\midrule
\endfirsthead
\multicolumn{4}{l}{\small (표 \thetable\ 계속)}\\
\toprule
\textbf{문헌} & \textbf{핵심 지표} & \textbf{조작화 방식} & \textbf{직전 연구 대비 심화 지점}\\
\midrule
\endhead
\bottomrule
\endfoot

\citet{handa-2025-which-economic-tasks-are} &
과업 사용 폭/깊이(breadth/depth) &
Clio로 대화를 O*NET 과업에 매핑, 직업별 "과업 대비 AI 사용 관측 비율"의 누적분포; 협업모드 5유형(자동화/증강) &
(AEI 원조) 최초로 "실사용"을 대규모 관측치로 정의\\

\citet{appel-2025-uneven-geographic-and-enterprise-ai} &
지리적 사용강도 &
AUI = 지역 사용점유율 / 지역 근로연령인구점유율; 1P API의 자동화비중(77\%) vs Claude.ai(50\%) &
Handa의 과업단위 사용을 지리$\cdot$소비자-기업 이원구조로 확장\\

\citet{tamkin-2025-estimating-ai-productivity-gains-from} &
과업당 시간절감률(depth) &
Claude 자신에게 "AI 없이/AI로 걸린 시간"을 추정시켜 $1-\text{time}_{\text{with}}/\text{time}_{\text{without}}$ 산출, Hulten 정리로 거시 생산성 집계 &
Handa의 "폭"(어떤 과업에 쓰이는가)을 넘어 과업 내 이질성(깊이)을 최초로 측정\\

\citet{appel-2026-economic-primitives} &
실효 AI 커버리지(성공률 가중) &
9개 신규 분류기(복잡성$\cdot$숙련$\cdot$용도$\cdot$자율성$\cdot$성공)로 "경제원시지표" 구축; 커버리지 = 성공률$\times$시간비중 가중합 &
단순 과업등장 여부(Handa)를 성공 여부로 할인해 "실질적" 영향으로 재정의\\

\citet{massenkoff-2026-labor-market-impacts-of-ai} &
관측노출(이론노출$\times$자동화 게이팅) &
$\beta\geq 0.5$ 과업만 인정 후 WorkUsage$\geq$100 임계치, 자동화비중 가중 $\alpha_t$로 직업단위 집계 &
Appel(2026)의 성공률 가중과 달리 이론적 상한($\beta$)을 게이트로 사용해 노출과 결합\\

\citet{massenkoff-2026-learning-curves} &
모델선택$\cdot$tenure(경력)별 사용 &
유료 이용자의 Opus 사용비중-임금 회귀; 고/저tenure 이용자 간 성공률 차이(3단계 회귀) &
횡단면 사용 스냅샷을 넘어 시간에 따른 사용자 학습(동태적 차원) 도입\\

\citet{fan-2026-what-countries-use-ai-and} &
국가별 사용 강도$\cdot$확산폭 &
강도=인구백만명당 대화수(PPML), 확산폭=25개 요청군집 로그HHI(OLS); AI-diffusion space(product space 응용) &
Appel(2025)의 지리적 AUI(국가 내 1개 스칼라)를 국가간 비교가능한 이중 지표+네트워크 구조로 확장\\

\citet{massenkoff-2026-cadences} &
산출물(artifact) 구성$\cdot$자율성$\cdot$보고/기대노출과의 연계 &
대화를 30여개 산출물로 분류, 토큰-임금 관계; AEI Survey(N$\approx$9,700)로 실사용 데이터와 주관적 인식을 프라이버시 보존 방식으로 연계 &
Massenkoff(2026c)의 tenure(행동 데이터 내부)를 넘어 설문이라는 외부 데이터원과 사용 로그를 사상 최초로 연계\\

\citet{fan-2026-aggregate-gains-from-ai} &
화폐화된 사용가치(LCE)$\cdot$분배(ACI) &
LCE=과업별시간절감$\times$임금$\times$사용비중의 국가합; ACI=보건경제학 concentration index를 임금분포-AI이득 분포에 응용 &
Tamkin$\cdot$McCrory(2025)의 미국 국내 시간절감 집계를 국제 비교+분배(누가 이득을 보는가)로 확장\\
\end{longtable}

이 표가 드러내는 계보는 뚜렷한 방향성을 갖는다: \textbf{폭(Handa) $\to$ 지리(Appel V3) $\to$ 깊이/시간가치(Tamkin$\cdot$McCrory) $\to$ 질/성공(Appel V4) $\to$ 동태성/학습(Massenkoff V5) $\to$ 국가간 확산(Fan) $\to$ 주관적 인식과의 연계(Massenkoff V6) $\to$ 화폐화$\cdot$분배(Fan$\cdot$Nguyen)}의 순서로, 매 단계가 "사용량"이라는 단일 개념 안에서 새로운 측정축(공간$\cdot$시간$\cdot$질$\cdot$주관성$\cdot$화폐가치)을 추가하는 방식으로 발전해 왔다.

% ============================================================
\section{논의의 중심축: 사용량 중심인가, 결합형인가}
\label{sec:axis}
% ============================================================

\begin{longtable}{L{3.2cm}L{2.8cm}L{9.0cm}}
\caption{9편의 중심 논의축 종합}\label{tab:axis}\\
\toprule
\textbf{문헌} & \textbf{중심축} & \textbf{근거}\\
\midrule
\endfirsthead
\multicolumn{3}{l}{\small (표 \thetable\ 계속)}\\
\toprule
\textbf{문헌} & \textbf{중심축} & \textbf{근거}\\
\midrule
\endhead
\bottomrule
\endfoot

\citet{handa-2025-which-economic-tasks-are} & 사용량 중심 & 노출도는 대조군으로만 언급, 핵심 기여는 실측 과업매핑 자체\\
\citet{appel-2025-uneven-geographic-and-enterprise-ai} & 사용량 중심 & 지리$\cdot$기업 차원의 사용 패턴 기술이 전부, 노출 개념 미등장\\
\citet{tamkin-2025-estimating-ai-productivity-gains-from} & 사용량 중심 & 시간절감이라는 사용량의 새 차원(깊이) 측정에 집중\\
\citet{appel-2026-economic-primitives} & 사용량 중심(성공률로 심화) & "실효 커버리지"가 노출과 유사한 개념이나 이론노출 미사용\\
\textbf{\citet{massenkoff-2026-labor-market-impacts-of-ai}} & \textbf{결합형(계산적 결합)} & Eloundou $\beta$(노출)를 실사용으로 게이팅해 단일 지수화, 노동시장 성과와 연결\\
\citet{massenkoff-2026-learning-curves} & 사용량 중심 & 모델선택$\cdot$tenure라는 순수 행동 데이터 분석\\
\citet{fan-2026-what-countries-use-ai-and} & 사용량 중심 & 노출지수를 서론에서 비판적으로 언급할 뿐 지표 자체는 순수 사용통계\\
\textbf{\citet{massenkoff-2026-cadences}} & \textbf{결합형(설문 연계)} & 설문 기반 주관적 노출인식과 관측사용 패턴(자동화비중)의 상관을 직접 추정\\
\citet{fan-2026-aggregate-gains-from-ai} & 사용량 중심(화폐화) & 노출지수를 비판적 서론 대비축으로만 사용, LCE$\cdot$ACI는 순수 사용 기반\\
\end{longtable}

종합하면, AEI 계열 9편의 \textbf{7편(78\%)은 순수 사용량 중심} 연구다. 이들에게 "노출도"는 독립적으로 측정$\cdot$투입되는 변수가 아니라, 서론에서 "우리가 넘어서는 선행 패러다임"으로 소환되는 수사적 장치에 가깝다—\citet{handa-2025-which-economic-tasks-are}의 "이론적 노출지수와 대비되는 실측 접근", \citet{fan-2026-what-countries-use-ai-and}$\cdot$\citet{fan-2026-aggregate-gains-from-ai}의 "잠재력이 아닌 실현"이라는 반복되는 문구가 이를 뒷받침한다. 반면 \textbf{2편(22\%)만이 노출도를 계산적으로 사용량과 결합}하며, 그 결합 방식도 서로 다르다—\citet{massenkoff-2026-labor-market-impacts-of-ai}는 \textbf{객관적$\cdot$이론적 노출(LLM 능력 평가)}을 실사용의 게이트로 삼고, \citet{massenkoff-2026-cadences}은 \textbf{주관적$\cdot$설문 기반 노출(사용자 자신의 인식)}을 실사용 패턴과 상관분석한다. 이 두 결합 방식은 각각 본 프로젝트의 \texttt{data/proc/merged/} 설계에 참고할 두 가지 서로 다른 템플릿을 제공한다(\S\ref{sec:conclusion}).

% ============================================================
\section{이전 연구의 구체적 계승$\cdot$발전 경로}
\label{sec:lineage}
% ============================================================

\subsection{노출도 연구의 계승: Massenkoff and McCrory(2026)의 "관측노출"}

\citet{massenkoff-2026-labor-market-impacts-of-ai}는 \citet{eloundou-2023-gpts-are-gpts-an-early}의 과업별 이론노출등급 $\beta$를 그대로 가져와, 과업 $t$의 관측노출을 다음과 같이 정의한다:
\[
\tilde{r}_t = \mathbf{1}\{WorkUsage_t \geq 100\} \times \mathbf{1}\{\beta_t \geq 0.5\} \times \alpha_t,
\]
여기서 $\alpha_t = 1/2 + 1/2 \times AutomationShare_t$는 자동화 사용비중에 따른 가중치다. 이를 O*NET 과업의 시간비중으로 가중평균해 직업단위 지수 $R_o$를 산출한다. 이 방식의 핵심은 \textbf{$\beta$가 "상한"으로만 기능}한다는 점이다—$\beta=1$(LLM 단독으로 충분)인 과업이라도 실제 WorkUsage가 임계치 미만이면 노출 0으로 처리된다. 그 결과, Computer \& Math 직업군은 이론상 94\%가 노출 가능하지만 실제 관측 커버리지는 33\%에 불과함을 보여, 이론적 잠재력과 실현 사이의 괴리를 정량화한다. 이는 \citet{eloundou-2023-gpts-are-gpts-an-early}이 제시한 "가능성"을 부정하지 않으면서도, 그 가능성 중 \textbf{실제로 실현된 부분만}을 걸러내는 방식으로 선행 노출연구를 계승$\cdot$발전시킨 사례다.

주목할 점은 이 논문이 계승하는 대상이 \citet{felten-2021-occupational-industry-and-geographic-exposure}의 AIOE나 \citet{webb-2020-impact-artificial-intelligence-labor}의 특허기반 지수가 \textbf{아니라} \citet{eloundou-2023-gpts-are-gpts-an-early}의 LLM 평가 기반 지수라는 것이다. 이는 \texttt{02\_AI노출도\_지표\_비교.tex}가 정리한 "LLM 영향의 조작화" 4단계(비특정 $\to$ LLM특화 $\to$ LLM평가자화 $\to$ 이론$\times$실사용 결합)에서 AEI 계열이 정확히 4단계(이론$\times$실사용 결합)에 위치하며, 그 이전 단계인 Eloundou의 3단계(LLM평가자화) 산출물을 직접 입력으로 삼는다는 계보를 보여준다. 즉 AEI 계열은 Felten$\cdot$Webb류의 "인간 전문가 판정" 노출지수 계보와는 단절되어 있고, "LLM 자신이 판정하는" 노출지수 계보(Eloundou)와만 계산적으로 연결된다.

\subsection{사용량/채택 연구의 계승: 자기보고에서 관측 로그로}

AEI 계열이 사용량 측정에서 계승$\cdot$발전시킨 선행연구는 노출도 문헌이 아니라 오히려 \textbf{경제학 일반이론과 채택 측정 전통} 두 갈래다.

첫째, \citet{fan-2026-what-countries-use-ai-and}는 국가별 사용을 강도(intensity)와 확산폭(breadth) 두 축으로 분해하는데, 이는 일반목적기술(General Purpose Technology, Bresnahan and Trajtenberg 1995)이 "얼마나 많이 쓰이는가"와 "얼마나 여러 부문에 퍼지는가"로 구분되어야 한다는 경제사적 이론을 계승한 것이다. 동시에 \citet{fan-2026-what-countries-use-ai-and}는 이 두 축을 측정하는 데 전통적 방식이었던 \textbf{기업 서베이}(Calvino et al. 2022; Calvino and Fontanelli 2023)를 명시적으로 대체하는 실측 대안으로 자신의 연구를 위치시킨다—"직업-과업 노출지수나 기업 서베이가 아니라 실제 Claude AI 대화 데이터를 이용"했다는 문구가 이를 직접 보여준다.

둘째, \citet{appel-2025-uneven-geographic-and-enterprise-ai}는 자신의 AUI(지리적 사용지수)가 그려내는 채택 그림을 미국 Census Business Trends and Outlook Survey(BTOS)의 자기보고 도입률(2023년 가을 3.7\%$\to$2025년 8월 9.7\%)과 나란히 인용하며 대비시킨다. 이는 채택 서베이 문헌의 결과를 부정하지 않으면서도, 서베이가 포착하지 못하는 \textbf{채택 이후의 활용 강도$\cdot$패턴}을 AEI가 보완할 수 있음을 시사하는 방식의 계승이다.

셋째, \citet{massenkoff-2026-cadences}의 AEI Survey는 가장 직접적인 방법론적 융합 사례다. 이 설문은 응답자의 자기보고(보고노출$\cdot$기대노출, 향후 고용전망 인식 등)를 수집한다는 점에서 \citet{bick-2026-mind-the-gap-ai-adoption}류의 근로자 서베이 전통을 그대로 잇지만, 결정적으로 \textbf{응답자의 실제 Claude 사용 세션(최대 20개)을 프라이버시 보존 방식으로 직접 연계}한다. 이는 자기보고 서베이 전통과 실사용 로그 전통이라는, 지금까지 별개로 발전해 온 두 계보를 개별 응답자 수준에서 사상 처음으로 결합한 사례로—채택 문헌의 "자기평가 편향" 한계와 사용 로그 문헌의 "동기$\cdot$인식을 관측할 수 없다"는 한계를 동시에 보완하려는 시도다.

\subsection{AEI 데이터의 방법론적 타당성에 대한 비판적 검증: Bick et al.(2026)와의 접점}
\label{sec:bick}

AEI를 1차 자료로 사용하지는 않지만, AEI가 산출한 사용량 지표 자체의 타당성을 외부에서 검증한 연구로 \citet{bick-2026-what-work-does-generative}가 있다. 이들은 Real-Time Population Survey(RPS)라는 별도의 미국 근로자 서베이로 O*NET 과업단위 genAI 채택지수를 구축한 뒤, \citet{handa-2025-which-economic-tasks-are}의 AEI 과업점유율 데이터와 비교해 상관계수가 0.34에 불과함을 보인다—AEI 등 챗로그 기반 지표는 지니계수(0.83)가 서베이 기반 지표(0.64)보다 훨씬 높아, \textbf{소수 과업(주로 문서편집 등 범용 활동)에 과대분류되는 경향}이 있음을 지적한다. 이는 AEI 계열이 스스로는 "채택 서베이보다 우월한 실측"이라 주장하는 근거(\S\ref{sec:usage})에 대해, 반대 방향의 자료(서베이)를 쓰는 연구자가 "챗로그 분류기 역시 편향에서 자유롭지 않다"는 반론을 제기하는 구도로, 향후 본 프로젝트가 AEI 데이터를 채택할 때 이 상호 비판을 함께 고려해야 함을 시사한다.

% ============================================================
\section{결론: 본 프로젝트에 대한 시사점}
\label{sec:conclusion}
% ============================================================

AEI 계열 9편을 종합하면, 이 연구 프로그램의 정체성은 \textbf{"노출도 연구를 대체하는 사용량 연구"}에 가깝다—9편 중 7편이 순수 사용량 중심이며, 노출도는 서론에서 소환되는 대비축일 뿐 변수로 결합되지 않는다. 그러나 나머지 2편은 본 프로젝트의 목표와 정확히 같은 문제의식—노출만으로는 포착 못하는 실제 채택$\cdot$활용 효과—을 이미 실현하고 있으며, 서로 다른 두 결합 방식을 제시한다.

\citet{massenkoff-2026-labor-market-impacts-of-ai}의 "이론노출$\times$실사용 게이팅" 방식은 기존에 이미 계산된 노출지표(예: 한국형 AIOE-KR, 또는 \texttt{02\_AI노출도\_지표\_비교.tex}가 정리한 8개 지표 중 하나)가 있을 때, 이를 상한으로 삼아 실제 사용 데이터로 "얼마나 실현되었는가"를 걸러내는 방식으로 즉시 응용 가능하다. 반면 \citet{massenkoff-2026-cadences}의 "설문-사용로그 연계" 방식은 본 프로젝트가 향후 자체 서베이를 설계할 경우, 자기보고 인식과 실제 사용 데이터를 개별 응답자 수준에서 연결하는 설계의 모범사례가 된다. 두 방식 모두 \texttt{data/proc/merged/}가 지향해야 할 "노출$\times$사용 결합" 데이터셋 설계 시, 매칭 키를 (i) 이론노출 산출 시 사용된 직업$\cdot$과업 코드, (ii) 실사용 데이터의 O*NET/SOC 코드, (iii) 설문 응답자 식별자(사용 시) 세 층위로 명시해 두어야 함을 시사한다.

% ============================================================
\printbibliography[title={참고문헌}]
% ============================================================

\end{document}
